\documentclass[12pt]{article}
\usepackage[utf8]{inputenc}
\usepackage[T1]{fontenc}
\usepackage[margin=1in]{geometry}
\usepackage{setspace}
\usepackage{xcolor}
\usepackage[colorlinks=true, urlcolor=blue!60!black, linkcolor=blue!60!black]{hyperref}
\setlength{\parindent}{0em}
\setlength{\parskip}{0.45em}
\renewcommand{\baselinestretch}{1.0}
\begin{document}
\begin{center}
{\Large \textit{Milliken v.\ Bradley}, 418 U.S.\ 717 (1974)\par}
\vspace{0.3em}
{\normalsize Case Brief --- One-Page Edition\par}
\end{center}

\textbf{Memory Jogger.} No metropolitan school desegregation remedy across suburban lines without an inter-district violation (\href{https://supreme.justia.com/cases/federal/us/418/717/}{Justia}).

\textbf{Facts.} Detroit's schools, about 63.6\% Black, were found \textit{de jure} segregated by the Detroit Board and the State. No proof was taken that any of the suburban districts had segregated or that their lines were racially drawn.

\textbf{Procedural History.} The district court ordered a remedy covering Detroit and 53 suburban districts; the Sixth Circuit affirmed.

\textbf{Votes.} Reversed, 5--4: Burger for the Court; Stewart concurred; Douglas, White, and Marshall dissented.

\textbf{Holding.} The violation bounds the remedy: cross-district relief requires that discriminatory acts substantially caused inter-district segregation, or that lines were drawn by race.

\textbf{Reasoning.} The Constitution does not require metropolitan racial balance; district lines reflect deep traditions of local control, and no cross-district effect was shown. Stewart would allow relief for racially redrawn lines or discriminatory housing laws. The dissenters said the State cannot hide behind its subdivisions; Marshall called it ``a giant step backwards.''

\textbf{Examples.} \emph{Same-Side:} gerrymandering inside one city. \emph{Opposite-Side:} a white town's secession to form its own district. \emph{Fence-Sitter:} suburban racial zoning and redlining.

\textbf{Critique.} Defenders praise its limits on federal equity power; critics, including Myron and Gary Orfield, say it ignored state-sponsored housing segregation and, with \textit{Rodriguez}, sheltered white flight.

\textbf{Key Quotation.} ``No single tradition in public education is more deeply rooted than local control over the operation of schools; local autonomy has long been thought essential both to the maintenance of community concern and support for public schools and to quality of the educational process.'' --- Chief Justice Burger, 418 U.S.\ at 741--42.

\end{document}
